\begin{proof}
Se trata de demostrar, teniendo en cuenta (6.13.2), que el punto
genérico de $\operatorname{Spec}(A')$ es interior a
$\operatorname{Nor}(\operatorname{Spec}(A'))$. La demostración sigue
el mismo camino que la de (6.12.4.1), cuyas notaciones conservamos.
Se observa primero que se puede suponer que $A$ es integralmente
cerrado, y entonces se sabe que $A_1=A[t_1,\ldots,t_n]$ está
integralmente cerrado (Bourbaki, \emph{Alg. comm.}, cap. V, § 1, n\textsuperscript{o} 3,
cor. 2 de la prop. 13); después la cuestión se reduce al caso en que $A'$ es
un $A$-módulo \emph{libre de tipo finito}; el razonamiento de (6.12.4.1)
demuestra entonces (reemplazando si hace falta $A$ por un anillo
$A_f$ con $f\neq0$) que las fibras $g^{-1}(z)$ del morfismo
$g:X'\to X$ son regulares y, \emph{con mayor razón}, normales. Además,
$g$ es plano y $X$ es normal, luego, por (6.5.4, (ii)), $X'$ es normal.
\end{proof}
